\fancypagestyle{normal}
\fancyhead{} % clear all header fields 
\fancyhead[RO]{Humanités, Littérature, Philosophie -- Banque des sujets indexés}
\fancyhead[LE]{\rightmark}
\fancyfoot{} %clear all footer fields
\fancyfoot[RO]{\textbf{\thepage} \hspace*{1.5cm}}
\fancyfoot[LE]{\hspace*{1.5cm}\textbf{\thepage}}
\renewcommand{\headrulewidth}{0.4pt}
\renewcommand{\footrulewidth}{0pt}

{\let\clearpage\relax%
\noindent\chapter{Avertissement}
}
\markright{Avertissement}{}
\label{avertissement}

Ce document vise à être un outil efficace pour l'ensemble des professeurs chargés d'enseigner
la spécialité «~Humanités, Littérature, Philosophie~». Il ne revêt aucun caractère officiel
et a été d'abord conçu pour simplifier le travail des professeurs, notamment par présence de différents index à la fin de ce document (à partir de la page \pageref{indexdeb}).

Si vous souhaitez accéder à une vision synoptique des différentes questions qui ont été posées, ainsi qu'aux sujets téléchargeables en version Word (.doc ou .docx) ou LibreOffice (.odt), vous pouvez vous rendre sur le Google Doc des \href{https://docs.google.com/document/d/1594BtSBMpnk_i"Lpqf7vzIPnFkFG6ZFvygOAMzdtbPhA/edit?usp=sharing}{Annales de sujets d'Audrey Pardigon}. Si vous souhaitez seulement avoir un fichier avec tous les sujets et un index des auteurs, vous les trouverez à \href{https://drive.google.com/file/d/1DnEUzL6_MTiRfS60gFycgNy-jxfQFMsu/view?usp=sharing}{ce lien} (odt) ou \href{https://docs.google.com/document/d/1XnJ397bNF5Vn-ZwEVjEvbJiCvbMBvRGi/edit?usp=sharing&ouid=103662385121242831300&rtpof=true&sd=true}{ce lien} (docx) (sans l'index des auteurs pour le fichier Word). Vous pouvez retrouver ce fichier et les fichiers qui ont servi à créer ce fichier (écrit en \LaTeX{}) \href{https://drive.google.com/drive/folders/1QwtvLY3i87WzyajRLPe0Xmmb31FK9oDp?usp=sharing}{ici}. Si vous souhaitez disposer d'une version de ce fichier avec les éléments d'évaluation, vous la trouverez ici~: \sujetshlpcorr{}.

\section*{Intégration des «~sujets zéro~» et ordre des sujets}
\phantomsection\addcontentsline{toc}{section}{«~Intégration des sujets zéro~» et ordre des sujets}

Quoiqu'ils n'aient pas fait l'objet d'une session d'examen, et qu'aucun élève n'ait donc composé sur ces sujets, on a intégré dans ces annales les «~sujets zéro ~», \cad{} les sujets proposés aux professeurs pour qu'ils se fassent une idée des attendus de l'examen. Ils ont été placés en premier pour chacun des grands thèmes de l'année de Terminale (la recherche de soi et l'humanité en question) et sont signalés, tant pour les index que pour le sujet lui-même, par le fait qu'on a indiqué «~Sujet 0 ~» comme centre d'examen. Ils peuvent ainsi être aisément distingués des autres sujets. Ce sont également les seuls sujets qui aient été publiés en 2020.

Pour l'ordre des sujets, on a procédé comme on avait déjà fait pour les \href{https://drive.google.com/file/d/1KN7gwAOkmIXFNAGYKYemD3y7Q0n68Phn/view?usp=sharing}{sujets de Première} en séparant les deux thèmes de l'année, de manière à permettre de naviguer plus aisément entre des sujets qui portent sur des thématiques similaires. On a pour chaque thématique classé les sujets par ordre chronologique d'année puis par ordre alphabétique du centre d'examen (Amérique du Nord, Amérique du Sud, Asie-Pacifique, Centres étrangers, La Réunion, Liban, Métropole, Nouvelle-Calédonie, Polynésie). Certains sujets n'ont parfois pas pu être retrouvés (à l'image des sujets du jour 1 pour les épreuves de spécialité du bac 2021 en Asie) ou n'ont jamais été publiés parce que les épreuves ont été totalement annulées, que ce soit en raison de l'épidémie de Covid-19 lors des sessions 2021 et 2022 ou pour des raisons politiques (les épreuves de spécialités ont ainsi été annulées en Nouvelle-Calédonie en 2021, et pour des raisons liées à la situation politique en 2024). Depuis 2022, les épreuves ayant lieu en-dehors de la Métropole sont de plus en plus regroupées par «~groupe ~» de Centres étrangers~: les épreuves de spécialités ont par exemple ainsi eu lieu au même moment en 2022 dans les Académies de Guadeloupe, Guyane, Mayotte, Martinique, Métropole et de la Réunion, ce qui explique l'absence de sujets spécifiques à la Réunion pour la session 2022. C'est aussi le cas du Liban en 2024 qui a rejoint les «~Centres étrangers gr.~1~».

\section*{Au sujet des index}
\phantomsection\addcontentsline{toc}{section}{Au sujet des index}

On trouvera à la fin, à partir de la page \pageref{indexdeb}, un \hyperlink{indexa}{index des auteurs} sur lesquels portent les sujets (p.~\pageref{indexa}), \hyperlink{indexg}{un index des textes selon leur genre} (p.~\pageref{indexg}), 
un \hyperlink{indexw}{index des questions selon les axes du programme puis selon leur type} (p.~\pageref{indexw}), ainsi qu'\hyperlink{indexv}{un index des questions selon leur type puis selon les thèmes du programme} (p.~\pageref{indexv}), \hyperlink{indexq}{un index général des questions selon leur type (interprétation ou réflexion et littéraire ou philosophique)} (p.~\pageref{indexq}), un \hyperlink{indexn}{index des sujets par lieux et années} (p.~\pageref{indexn}), qui permet de se faire une idée rapide des sujets proposés au fil des années et de retrouver l'ordre dans lequel les sujets ont été publiés,  ainsi qu'\hyperlink{indext}{un index des thèmes abordés par les textes}, qui permet de chercher rapidement un sujet sur un thème abordé en cours (p.~\pageref{indext}). Pour se repérer plus rapidement dans cet index, nous renvoyons à la \hyperlink{repere-thematique}{section proposant une liste de termes indexés} pour chaque chapitre et grandes thématiques au sein de chaque chapitre (p.~\pageref{repere-thematique}).

Plusieurs des index (concernant les axes du programme particulièrement, et concernant les genres des textes)
sont une simple proposition et n'engagent que moi.
J'ai essayé au mieux d'identifier le genre du texte, ce qui a parfois été difficile,
particulièrement pour les textes littéraires, dont je ne suis pas spécialiste.
Doit-on bien ranger l'ouvrage de Simone de Beauvoir, \textit{Mémoires d'une jeune fille rangée} dans la catégorie des autobiographies, quand cela a également une portée philosophique et qu'à certains égards on pourrait le considérer comme une forme d'apologie~? Que dire de l'ouvrage de Charlotte Delbo qui, certes, est le récit de l'expérience concentrationnaire de l'auteur, mais qui emploie une multitude de genres pour le faire~? Par nécessité j'ai tranché, mais cette catégorisation n'est pas absolue, et je serais heureux 
d'\href{mailto:mikael.quesseveur@ac-lille.fr}{avoir des retours}
et d'être corrigé.

On a essayé, de notre mieux, de rattacher chaque sujet aux grands axes de chaque semestre, la distinction entre les «~Expressions de la sensibilité~» et les «~Métamorphoses du moi~» étant néanmoins parfois poreuses, tout comme le rapport entre violence et technique qui rattache les «~Limites de l'humain~» à «~Histoire et violence~». Parfois, la classification est originale. Certains mêlent ainsi les deux semestres, à l'image du sujet d'essai littéraire sur le texte d'Hans Jonas (p.~\pageref{jonas-amérique2021}) sur la technique, qui interroge la capacité de la littérature à nous rendre sensible à l'importance de la nature, thématique qu'on aurait plutôt attendu autour des expressions de la sensibilité, mais qui n'a pris son sens qu'en considérant la question de la technique et de l'écologie. D'autres sujets semblent exiger des connaissances issues des cours de Première, à l'image des sujets portant sur un passage de \textit{Tristes tropiques} de Claude Lévi-Strauss (p.~\pageref{lévi-strauss-liban2022}) qui fait appel pour l'interprétation philosophique au concept de monde et à la question de la découverte, et pour l'essai littéraire à la question de la culture. Ce n'est qu'en tant que ce sujet se rapporte au concept d'humanité qu'il semble bien pouvoir se rapprocher explicitement de «~L'humanité en question ~», et que l'on peut certainement y voir une interrogation sur la modernité qui, à certains égards, introduit une pensée post-moderne qu'on peut rattacher aux «~Limites de l'humain ~».

Pour autant, cette catégorisation peut certainement être remise en question et je serais heureux d'\href{mailto:mikael.quesseveur@ac-lille.fr}{avoir des retours} et d'être corrigé. On remarquera d'ailleurs, qu'entre la première version et celle-ci, j'ai effectué certaines corrections, que l'on retrouvera dans les 
\hyperlink{errata}{Errata et corrections}.

\section*{Les statistiques}
\phantomsection\addcontentsline{toc}{section}{Les statistiques}

Puisqu'il commence à y avoir beaucoup de sujets, on a également, à partir du travail effectué, ajouté un chapitre de \hyperlink{statistiques}{Statistiques}, qui permet de se faire une idée rapide des sujets qui ont jusqu'à présent été proposés, de manière à en évaluer la répartition thématique notamment et les éventuelles affinités entre les types de sujets et les types de questions. Il convient de rester néanmoins très prudent~: «~Éducation, transmission, émancipation~»{} et «~Création, continuités et ruptures~»{} n'ayant pu faire l'objet d'une évaluation qu'à partir de la session 2024, il existe un écart important entre le nombre de sujets proposés sur ces deux chapitres et les autres~; l'écart est d'autant plus important que pendant la période du Covid, chaque jour étaient proposés deux sujets, et que «~Éducation, transmission, émancipation~»{} ni «~Création, continuités et ruptures~»{} ne pouvaient faire l'objet d'un sujet à ce moment.

Ces statistiques sont divisées en deux grands groupes :

\begin{enumerate}
\item D'abord les statistiques générales sur l'ensemble des sujets. On a ainsi détaillé selon les disciplines pour chaque chapitre le nombre de questions posées, en essayant de rendre l'ensemble le plus lisible possible. Pour rendre l'ensemble le plus efficace possible, les nombres dans le tableau renvoient désormais directement par un hyperlien à \hyperlink{indexv}{l'index selon les types de sujet puis selon les selon les axes du programme}, ce qui permet rapidement de faire la liste des questions posées pour chaque thème. Un troisième tableau résume l'ensemble des questions posées selon les thèmes, de manière plus compacte, et sans répéter les données\footnote{On a conservé les trois tableaux parce qu'il paraît utile pour un professeur, qui cherche à identifier le nombre de questions par axes du programme de manière précise de suivre la logique des chapitres du programme, en cherchant dans l'entrée pertinente du programme, ce que ne permet pas le tableau plus compacte, qui demande peut-être plus d'habitude pour être lu efficacement.} On a également inséré, à la suite une page qui résume, à l'aide de graphiques, la répartition des sujets selon les axes du programme. Ce dernier tableau et cette page de graphique sont repris dans les pages qui suivent pour les statistiques par centres d'examen.

\item Ensuite les statistiques par centres d'examen, classés par ordre alphabétique. Pour que ces statistiques soient suffisamment significatives, on a intégré que les centres d'examen dans lesquels il y a eu suffisamment de sujets proposés, pour que cela ait du sens à représenter graphiquement les données. On a essayé de rendre ces graphiques aussi lisibles que possibles, mais comme pour les tableaux précédents, il était nécessaire de faire appel aux acronymes que l'on utilise ailleurs. Il permet certainement en tout cas de se faire une idée rapide des tendances (passées) dans les différents centres d'examen.
\end{enumerate}

\section*{Corrections apportées aux sujets}
\phantomsection\addcontentsline{toc}{section}{Corrections apportées aux sujets}

Pour certains sujets, on s'est permis d'effectuer des corrections d'ordre typographique. Ainsi a-t-on pu mettre \textit{La Marseillaise} en italique dans le sujet portant sur \textit{14} de Jean Echenoz (Polynésie, 2021) ou encore mettre les majuscules qui conviennent à «~Première Guerre mondiale ~», conformément aux normes, dans l'intitulé du sujet d'interprétation philosophique portant sur le texte de Walter Benjamin, \textit{Expérience et pauvreté} (Amérique 2022). On a corrigé également une coquille manifeste dans le «~sujets 0 ~» portant sur une lettre de Jacques Rivière à Antonin Artaud. Dans le premier paragraphe, en effet, on a supprimé des guillemets qui étaient répétés deux fois~: \textcolor{darkgray}{«~qui du dehors attaque l’âme~»~»} est ainsi devenu \textcolor{darkgray}{«~qui du dehors attaque l’âme~»}. On a également modifié, selon les règles typographiques traditionnelles du français, la présentation des ordinaux de siècle, en plaçant en exposant l'abréviation du suffixe des adjectifs ordinaux, notamment pour le sujet portant sur Michel Foucault proposé en Métropole le jour 1 en 2021, ou encore dans les notes de bas de page du sujet proposé le deuxième jour en Polynésie en 2024. En-dehors de ces corrections typographiques et de la correction de la coquille, on n'a corrigé aucun texte des sujets. Pour les autres corrections, on pourra voir les corrections apportées à la précédente édition (notamment des coquilles) dans les \hyperlink{errata}{Errata et corrections}.

\section*{Remarques concernant les sujets et leur indexation}
\phantomsection\addcontentsline{toc}{section}{Remarques concernant les sujets}

Qu'on me permette ici d'insister sur le fait que l'indexation que je propose pour ces sujets n'a aucune valeur absolue ni officielle. Elle est très certainement contestable en bien des moments, et si j'ai essayé, au mieux, d'identifier le chapitre ou les chapitres de référence pour chaque question, il est possible que certains de mes choix ne soient pas les plus adéquats. N'hésitez pas à ce titre, si vous trouviez l'indexation de certains sujets contestables, à \href{mailto:mikael.quesseveur@ac-lille.fr}{me l'indiquer} pour que j'effectue les corrections nécessaires. Pour cette nouvelle édition, j'ai d'ailleurs moi-même changé l'indexation de certains sujets (ce qui est signalé dans \hyperlink{errata}{les \textit{errata} et corrections}.

À ce titre, il convient de rester prudent concernant les statistiques ou même le classement de ces sujets au sein de deux chapitres différents correspondant aux axes du programme~: ils sont l'effet de mes choix et n'ont aucune valeur en soi. Les statistiques, particulièrement, ne dépendent que du fait que j'ai indexé les sujets d'une certaine manière, et on obtiendrait certainement des résultats très différents avec une indexation différente. Concernant les sujets qui font appel à des connaissances des deux semestres, lorsqu'un des deux questions se référaient explicitement à un seul semestre, on a placé le sujet dans le chapitre correspondant. Quand cela n'était pas le cas, on a essayé d'identifier la thématique majeure du texte, et on a ainsi placé le sujet là où cela paraissait pertinent. Si les thématiques étaient trop intrinsèquement mêlées, on a généralement placé le sujet dans «~La recherche de soi~~». 

Puisque le nombre de sujets commence à devenir plus conséquent, il devient possible de faire quelques remarques qui permettent, peut-être, de mieux appréhender l'épreuve. On remarquera à ce titre quelques singularités. La première, et peut-être la plus instructive (et surprenante) réside dans le fait que deux textes identiques ont été proposés en 2023 en Nouvelle-Calédonie (Jour 1) et en Métropole (session de Remplacement, Jour 2), à savoir un texte de Joseph Ponthus extrait d'\textit{À la ligne, Feuillets d'usine}. Les deux sujets ne sont pas strictement identiques, puisqu'on ne trouve au v.~19 dans le sujet de Nouvelle-Calédonie aucune espace qui précède «~Tous les matins du monde ~», tandis que c'est le cas pour le sujet proposé en Métropole (v.~22 dans le sujet officiel). La numérotation des lignes y est également différente, puisque les marges semblent avoir été plus grandes pour le sujet de Métropole. Nous avons, pour la mise en page en général des vers, privilégié celle de Nouvelle-Calédonie, puisqu'il semblait alors que les retours à la ligne étaient moins dus à la disposition spécifique du sujet qu'à une reproduction du sujet (chaque vers commençant ainsi par une majuscule). Néanmoins, on a reproduit la différence entre les deux sujets quant à la présence de l'espacement. On remarquera également que les questions proposées n'étaient pas identiques, puisque le sujet calédonien propose comme question d'interprétation~: «~Comment s'exprime, dans ce texte, la violence de l'usine sur «~toutes les ouvrières et tous les ouvriers~~»~?~~», tandis que le sujet métropolitain propose~: «~«~Je suis là sans y être~~», dit le poète. En quoi cette phrase rend-elle compte du travail à l'usine tel qu'il est décrit dans ce texte~?~~». Pour les sujets de réflexion nous avons en Nouvelle-Calédonie~: «~Peut-on rester humain dans des conditions inhumaines~?~»  et en Métropole~: «~Le travail menace-t-il notre dignité~?~~». C'est l'occasion ainsi de comprendre que les questions posées ne sont pas rigidement attachées à chaque texte et qu'il est possible pour tout professeur de choisir une question qui lui paraîtrait pertinente pour préparer ses élèves au regard de son cours et de ce qu'il trouve intéressant dans un texte.

On trouve un phénomène relativement identique en Polynésie en 2022 avec un texte d'Hannah Arendt quasiment identique au sujet~0 proposé sur Hannah Arendt. En fait, seule la première phrase du texte du sujet~0 a été enlevée, et le reste du texte est commun. Les questions à nouveau sont différentes, et abordent, pour l'essai autant que pour la question d'interprétation, des dimensions radicalement différentes du texte ou de la pensée. Il est en ce sens envisageable que des sujets déjà proposés lors de sessions précédentes puissent faire, à nouveau, l'objet d'une épreuve avec des questions différentes.

\section*{Éléments d'évaluation}
\phantomsection\addcontentsline{toc}{section}{Éléments d'évaluation}

Cette version du fichier ne contient pas les éléments d'évaluation. Si vous souhaitez accéder à la version de ce fichier avec les éléments d'évaluation, vous pouvez vous rendre à ce lien \sujetshlpcorr{}.

À partir de l'édition 2025, on a proposé une indexation des sujets disposant d'éléments d'évaluation. Ces éléments d'évaluation n'ont pas de valeur prescriptive, mais permettent certainement d'une part d'identifier des références pertinentes pour les essais (et donc d'enrichir le cours) et d'autre part pour les questions d'interprétation de s'assurer que l'on a correctement identifié les éléments saillants d'un texte quant à son interprétation. Ils visent à être un guide pour les professeurs dans la correction des épreuves, et il nous paraissait en ce sens pertinent de les intégrer dans une version au moins de la banque indexée des sujets. On a ajouté, en plus de ces éléments d'évaluation un index des références culturelles mobilisées dans les éléments d'évaluation. L'ajout des éléments d'évaluation a le principal inconvénient d'allonger (très) largement le document : l'édition 2022 comptait au total 110~pages, tandis que l'édition comportant les éléments d'évaluation dépasse les 380~pages. Parce que \LaTeX{} est destiné à la manipulation de ce genre de documents, cela ne pose pas de difficultés dans la conception, mais il est possible à l'usage que cela ne soit pas idéal. Si vous préférez ce fichier à celui contentant les éléments d'évaluation, je suis intéressé par les raisons de vos préférences. N'hésitez pas, en ce sens, si vous utilisez cette banque de sujets, à \href{mailto:mikael.quesseveur@ac-lille.fr}{m'indiquer ce que vous en pensez}. Si vous souhaitez disposer d'un document avec les éléments d'évaluation vous le trouverez en cliquant sur ce lien~: \sujetshlpcorr{}.

\section*{Au sujet de ce PDF}
\phantomsection\addcontentsline{toc}{section}{Au sujet de ce PDF}

Le PDF que vous lisez a été fait sous \LaTeX{}, notamment en raison de la nécessité de disposer de plusieurs index
et des grandes possibilités offertes par ce langage pour l'automatisation d'un certain nombre de tâches
par l'utilisation des macros. Cela permettait, au demeurant, de rendre ce PDF «~interactif~», c'est-à-dire
que vous pouvez facilement naviguer au sein de ce PDF à l'aide des hyperliens cliquables qui s'y trouvent.
Vous trouverez également des signets (\textit{bookmarks}) qui vous permettent de naviguer facilement à partir du
panneau latéral de votre lecteur de PDF. Les liens sont indiqués par des textes avec une couleur différente~: en bleu les liens internes au document (présentés ainsi~: \textcolor{darkblue!80!black}{lien interne}) et en violet foncé les liens externes (présentés ainsi~: \textcolor{darkviolet!80!black}{lien externe}).
Les couleurs ont été choisies pour ne pas être (trop) visibles à l'impression tout en s'assurant qu'ils restent bien identifiables lors de la navigation dans le PDF.

On a conçu le document pour qu'il distingue les pages de gauche et de droite~: de la sorte on pourra aisément
l'imprimer en recto-verso.

Un mot enfin sur les droits de diffusion et de partage de ce fichier. Ce document et les fichiers \LaTeX{} qui
en sont à l'origine sont sous Licence Creative Commons : \doclicenseLongText{} Autrement dit vous êtes autorisés
à partager, distribuer et communiquer le matériel par tous moyens et sous tous formats ainsi qu'à adapter 
(c'est-à-dire à remixer, transformer, et créer) ce fichier à vos besoins pour toute utilisation non-commerciale.
Vous devez créditer l'Œuvre, intégrer un lien vers la licence et indiquer si des modifications ont été effectuées à
l'œuvre. Vous devez indiquer ces informations par tous les moyens raisonnables, sans toutefois suggérer que l'Offrant vous
soutient ou soutient la façon dont vous avez utilisé son œuvre. Vous n'êtes pas autorisé à faire un usage commercial de
cette œuvre, tout ou partie du matériel la composant. Dans le cas où vous effectuez un remix, que vous
transformez, ou créez à partir du matériel composant l'œuvre originale, vous devez diffuser l'œuvre modifiée dans les même
conditions, c'est à dire avec la même licence avec laquelle l'œuvre originale a été diffusée. Il va sans dire que les
fichiers de textes ne sont pas soumis à ces droits puisqu'ils appartiennent en partie au domaine public ou qu'ils sont des extraits d'œuvres sur lesquels je n'ai aucun droit. Par contre, l'utilisation du
code qui a permis la création de fichier est, lui, soumis à ces droits. À dire vrai ce n'est pas tant pour que mon
nom soit cité que je mets ce document sous Licence, mais pour éviter que mon travail ne soit commercialisé (qu'il
s'agisse du code ou de la première indexation des sujets que je propose). Pour autant, toute participation à 
l'amélioration de l'indexation est la bienvenue, et je vous engage à m'écrire à 
\href{mailto:mikael.quesseveur@ac-lille.fr}{cette adresse}.

\begin{flushright}M.~Q.\end{flushright}

\vfill{~}

\doclicenseImage